\subsecn{8}{D.8}{Exactly when join covers form a site}

\begin{thmplain}[\hypertarget{stmt:D.8}{Theorem~D.8}]

Let \(L\) be a finite lattice. Declare a nonempty family \((x_i\le x)_i\) to cover \(x\) when \(\bigvee_i x_i=x\). These covers are stable under pullback along \(y\le x\), namely \((x_i\wedge y)_i\) covers \(y\), iff \(L\) is distributive. In that case they form a subcanonical pretopology. The corresponding assertion is about the internal lattice operations of \(L\), not the ambient meet of a containing partition lattice.
\end{thmplain}

\begin{proof}

Distributivity gives \(\bigvee_i(x_i\wedge y)=(\bigvee_i x_i)\wedge y=x\wedge y=y\). Identity covers and composition of covers follow from singleton joins and associativity of joins. Conversely apply pullback stability to the cover \(\{a,b\}\) of \(a\vee b\) and the object \(y=c\wedge(a\vee b)\). It gives
\[c\wedge(a\vee b)=(a\wedge c)\vee(b\wedge c),\]
since \(a\wedge y=a\wedge c\) and likewise for \(b\). This distributive identity characterizes distributive lattices. For subcanonicity, the representable presheaf \(\operatorname{Hom}(-,z)\) on a poset is empty or a singleton at each object. Local sections on a cover mean \(x_i\le z\) for every \(i\), hence \(x=\bigvee_i x_i\le z\), which gives the unique amalgamation.
\end{proof}

\textbf{Counterexample and consequence.} In \(\mathrm{Part}(\{1,2,3\})\), put \(a=12|3\), \(b=13|2\), \(c=23|1\). The family \(\{a,b\}\) covers \(\top\), but its pullback along \(c\) is \(\{\bot,\bot\}\), which does not cover \(c\). Thus the general achievable-partition theory must retain the positive subset site of \hyperlink{sec:35}{§35}, restrict to a distributive achievable lattice, or specify another topology. It cannot simply rename all join equations a Grothendieck topology. This settles the direct join-cover question in P2; it does not identify arbitrary registration with sheafification.

\subsecn{9}{D.9}{When a witness defect changes the information partition}

\begin{thmplain}[\hypertarget{stmt:D.9}{Theorem~D.9}\ {\normalfont\upshape(cut criterion)}]

Let \(H=(U,E_H)\) and \(J=(U,E_J)\) be finite undirected graphs with \(E_J\subseteq E_H\). The component partition of \(J\) strictly refines that of \(H\) iff there is a nonempty proper subset \(A\) of some \(H\)-component \(C\) such that every \(H\)-edge between \(A\) and \(C\setminus A\) is absent from \(J\). Equivalently, \(J[C]\) is disconnected for some \(H\)-component \(C\). When exactly one edge is deleted, strict refinement occurs iff that edge is a bridge of \(H\).
\end{thmplain}

\begin{proof}

Each \(J\)-component lies within an \(H\)-component. If refinement is strict, choose one \(J\)-component \(A\) inside a split \(H\)-component \(C\). No \(J\)-edge leaves \(A\) inside \(C\), so every \(H\)-edge in that cut was deleted. Conversely such a cut prevents a \(J\)-path between its sides. Because \(C\) was \(H\)-connected, both sides belonged to the same old component; therefore the new partition is strictly finer. A single deleted edge can remove a nonempty cut exactly when it is a bridge.
\end{proof}

For §37.7 take \(H=\bigcap_v\Gamma_{\{v\}}\) and \(J=\Gamma_T\). This is an exact necessary-and-sufficient certificate for a nonzero witness \textbf{interval}. A nerve certifies the loss of a particular edge; D.9 supplies the missing step from that loss to an interval. Under a full-support prior, some question detects the strict interval by \hyperlink{stmt:53.2}{Proposition~53.2}; an arbitrary fixed question need not.

\textbf{Absorption example.} On six candidates with orbits \(\{1,2\},\{3,4\},\{5,6\}\), take \(v=(0,0,0,0,0,1)\) and \(w=(0,0,0,1,1,0)\). The singleton-graph intersection is a triangle and the joint graph removes only the edge between the first and third orbits. That pair's nerve consists of two isolated vertices and has nonzero reduced \(H_0\), but the remaining graph is connected. Its witness interval is zero, exactly as D.9 predicts.

\subsecn{10}{D.10}{Exact propagation on coordinate join trees}

\begin{thmplain}[\hypertarget{stmt:D.10}{Theorem~D.10}\ {\normalfont\upshape(finite constraint propagation)}]

Let the scene be a genuine finite Cartesian product \(D=\prod_{i\in I}D_i\). Let bags \(B_j\subseteq I\) cover \(I\) and form a tree with the running-intersection property: the bags containing any fixed coordinate form a connected subtree. At bag \(j\) let \(E_j\subseteq\prod_{i\in B_j}D_i\) be a constraint. For a directed tree edge \(j\to k\), let \(S_{jk}=B_j\cap B_k\) and define the message
\[M_{j\to k}=\operatorname{proj}_{S_{jk}}\left(E_j\cap\bigcap_{\ell\sim j,\ell\ne k}\operatorname{cyl}_{B_j}(M_{\ell\to j})\right),\]
computed from leaves inward on the component left after deleting \(jk\). Cylinders here restrict precisely the coordinates in the message's separator. Then
\[E_r\cap\bigcap_{j\sim r}\operatorname{cyl}_{B_r}(M_{j\to r})\]
is exactly the projection onto \(B_r\) of the global feasible set \(\bigcap_j\operatorname{cyl}_{I}(E_j)\). Empty constraints and empty separators are allowed: the latter have one assignment, and a message can be empty or that singleton.
\end{thmplain}

\begin{proof}

Induct over the component sending a message. The induction invariant is that the message comprises exactly the separator assignments that extend to a feasible assignment of that component's bags. A feasible component assignment gives the displayed projection by restriction. Conversely a bag assignment in the displayed intersection has, by induction, a feasible extension into each neighboring component. Running intersection implies that two such components share coordinates only through the central bag, and that a component shares coordinates with that bag only through its separator. The extensions therefore agree on every overlap and can be united. This proves the invariant and, at the root, the stated projection identity.
\end{proof}

This provides the exact local-computation result relevant to P7 on coordinate domains. It relies on the overlap/extension property supplied by the Cartesian scene and running intersection. Arbitrary partitions of a constrained common scene need not satisfy it; generalized saturation alone does not license this message-passing rule. The result concerns feasible constraints, not graded coupled channels or additive anomaly messages.

\subsecn{11}{D.11}{Finite certificates for polyhedral transfer classes}

\begin{thmplain}[\hypertarget{stmt:D.11}{Proposition~D.11}]

Let \(\mathcal A=\operatorname{conv}\{Q_1,\ldots,Q_N\}\) be a class of finite experiments on fixed state and output alphabets, and let \(\Gamma\) be a fixed finite budget experiment. Then
\[Q\preceq_B\Gamma\text{ for every }Q\in\mathcal A
\quad\Longleftrightarrow\quad
Q_j\preceq_B\Gamma\text{ for every }j.\]
More generally, \(\sup_{Q\in\mathcal A}\delta(\Gamma,Q)=\max_j\delta(\Gamma,Q_j)\).
\end{thmplain}

\begin{proof}

Necessity follows because each vertex belongs to the class. For sufficiency choose simulators \(G_j\) with \(Q_j=G_j\Gamma\). If \(Q=\sum_j\alpha_jQ_j\), the stochastic matrix \(G=\sum_j\alpha_jG_j\) gives \(Q=G\Gamma\). For the approximate statement use an optimal simulator for each vertex; convexity of TV gives \(\delta(\Gamma,Q)\le\sum_j\alpha_j\delta(\Gamma,Q_j)\le\max_j\delta(\Gamma,Q_j)\). The reverse inequality again uses membership of each vertex.
\end{proof}

For a single-record finite corruption model at fixed rate, linear normalization, nonnegativity, and prescribed knowledge-cell restrictions define a compact polytope of clutter kernels; the perceived channel is an affine image and hence a polytope. A declared finite policy family can similarly be convexified by candidate-independent randomization. D.11 reduces adequacy of a proposed budget to finitely many simulation LPs and supplies a maximizing attack vertex for the worst deficiency. This is a certificate for a \textbf{specified} transfer, not a proof of existence of a least transfer among all experiments. Arbitrary sequential factorizations and independently chosen products need not form this polytope without a separate realization argument.

\subsecn{12}{D.12}{Witness anomalies do not imply marginal contextuality}

\begin{thmplain}[\hypertarget{stmt:D.12}{Proposition~D.12}]

A coherent finite channel family \((P_T)_T\), restricted to any cover of a finite set \(T\), has a global distributional amalgamation at each candidate and under every prior. Consequently a nonzero admissibility or witness interval in such a family does not imply nonextendability of its marginal empirical model. Parity with its invariance ceiling supplies a nonzero witness interval in a model admitting a global distribution.
\end{thmplain}

\begin{proof}

At a candidate \(z\), the channel row \(P_T(\cdot\mid z)\) is the required global law by coherence. Under a prior \(\mu\), the mixture \(\sum_z\mu(z)P_T(\cdot\mid z)\) has all the specified mixed marginals because marginalization is linear. In the parity example, sample the two-bit candidate and report its coordinate values; this gives an explicit global law on both measurements, while §37.7 computes a nonzero witness interval under the orbit-invariance ceiling.
\end{proof}

The presheaf of candidate-pair witnesses uses inclusion restrictions retaining the entire pair. The event presheaf in contextuality uses coordinate restriction of outcome assignments; at the empty context it has one assignment, whereas \(W_\varnothing=O\times O'\) generally has several. Even their empty-context objects preclude an unqualified identification. An empty coupling fiber for merely local, mutually compatible channels is a genuine marginal-extension problem (§39); it cannot simultaneously be the restriction of a coherent global family. A richer construction relating selected witness obstructions to contextuality may be useful, but no such equivalence is needed for the finite witness calculus proved here.

\subsecn{13}{D.13}{Exactly when partition extractions commute}

\begin{thmplain}[\hypertarget{stmt:D.13}{Theorem~D.13}\ {\normalfont\upshape(rectangular intersection criterion)}]

Let \(a,b\) be partitions of a finite nonempty set \(D\), with saturation operators \(s_a,s_b\). The following are equivalent:

\begin{enumerate}
\def\labelenumi{\arabic{enumi}.}
\tightlist
\item
  \(s_a s_b=s_b s_a\) on every subset of \(D\).
\item
  \(s_a s_b=s_b s_a=s_{a\wedge b}\).
\item
  Within each block \(C\) of the ambient meet \(a\wedge b\), every \(a\)-block contained in \(C\) meets every \(b\)-block contained in \(C\).
\end{enumerate}

Equivalently, each connected component of the bipartite incidence graph of \(a\)-blocks and \(b\)-blocks is complete bipartite. This is a precise additional condition under which the two focusing operations commute.
\end{thmplain}

\begin{proof}

Write \(R_a,R_b\) for the equivalence relations of the partitions. Equality of the saturation composites on singleton sets is equality of the corresponding relational composites. If the composites commute, their common value is reflexive and symmetric; it is transitive because
\[(R_aR_b)(R_aR_b)=R_aR_aR_bR_b=R_aR_b.\]
It contains both relations and is contained in the equivalence relation they generate, so it equals \(R_{a\wedge b}\). This proves (1) implies (2); the reverse is immediate.
\end{proof}

Assume (2), choose an \(a\)-block \(A\) and a \(b\)-block \(B\) in the same meet block, and choose \(x\in B\). The set \(s_a s_b(\{x\})\) is the union of the \(a\)-blocks meeting \(B\). By (2) it is the whole meet block containing \(x\), so it contains \(A\), and \(A\cap B\ne\varnothing\). Conversely, (3) implies that starting in any \(b\)-block, one \(b\)-saturation followed by an \(a\)-saturation reaches its whole meet block; the reverse order does likewise. This proves (2) on singletons and hence on all subsets, because saturation preserves unions. Finally paths in the incidence graph are precisely alternating block-identification chains, whose unions are the meet blocks. Thus (3) is exactly the complete-bipartite condition. \(\square\)

The partitions \(12|3\) and \(1|23\) fail this condition: they have one meet block but their singleton blocks \(\{3\}\) and \(\{1\}\) do not meet. Coordinate projections of a Cartesian product satisfy it. The generalized set algebra of \hyperlink{sec:6}{§6} does not require this criterion for every label pair; invoking a commuting-extraction or propagation law requires checking the relevant intersection/extension conditions instead of inferring them from join closure.
